\documentclass[11pt]{article}
\usepackage[margin=1in]{geometry}
\usepackage{enumitem}
% =============================================================================
% Preambles/header.tex — shared LaTeX/Beamer preamble for this project
%
% Use from a Beamer lecture by prepending:
%     \input{header}
% and compiling with TEXINPUTS=../Preambles:$TEXINPUTS (the /compile-latex
% skill does this automatically).
%
% PALETTE CONTRACT. Color names defined here MUST match the names in
% Quarto/theme-template.scss so Beamer and Quarto renderings use the same
% palette. The scripts/check-palette-sync.sh script enforces this.
%
% To customize: edit the HEX values below AND the matching $variable in
% Quarto/theme-template.scss, then run scripts/check-palette-sync.sh.
% =============================================================================

% -----------------------------------------------------------------------------
% Engine detection + encoding
%
% Our /compile-latex skill uses XeLaTeX, where inputenc is unnecessary and
% can produce encoding warnings. Only load inputenc under pdfTeX.
% -----------------------------------------------------------------------------
\usepackage{iftex}
\ifPDFTeX
  \usepackage[utf8]{inputenc}
\fi

\usepackage{amsmath, amssymb, amsthm}
\usepackage{booktabs}
\usepackage{graphicx}

% -----------------------------------------------------------------------------
% PALETTE — mirrors Quarto/theme-template.scss variable names.
% Load xcolor and define colors BEFORE anything that references them
% (notably hyperref, hypersetup, and the Beamer color assignments).
% -----------------------------------------------------------------------------
\usepackage{xcolor}

\definecolor{primary-blue}{HTML}{012169}      % headings, accents
\definecolor{primary-gold}{HTML}{B9975B}      % emphasis, borders
\definecolor{highlight-yellow}{HTML}{F2A900}  % markers, alerts
\definecolor{light-bg}{HTML}{E8EDF5}          % title / section backgrounds
\definecolor{jet}{HTML}{1A1A1A}               % body text

% Semantic colors (used in diagrams and annotations)
\definecolor{positive}{HTML}{15803D}          % good / observed / identified
\definecolor{negative}{HTML}{B91C1C}          % bad / problematic / bias
\definecolor{neutral}{HTML}{525252}           % reference / context

% Highlight accents (match .hi-* classes in the SCSS)
\definecolor{hi-slate}{HTML}{314F4F}
\definecolor{hi-green}{HTML}{2E8B57}
\definecolor{hi-red}{HTML}{C41E3A}

% Semantic aliases so skill templates and snippets can write
% \textcolor{accent}{...} without hard-coding "primary-blue".
\colorlet{accent}{primary-blue}
\colorlet{accent2}{primary-gold}

% -----------------------------------------------------------------------------
% Hyperref — loaded AFTER xcolor + palette so link colors resolve.
% -----------------------------------------------------------------------------
\usepackage{hyperref}
\hypersetup{colorlinks=true, linkcolor=primary-blue, urlcolor=primary-blue,
            citecolor=primary-blue, pdfborder={0 0 0}}

% -----------------------------------------------------------------------------
% Course code — set once per deck (after \input{header}, before \begin{document})
% via \coursecode{CS401}. Shown in the footer on every slide so a deck's course
% is identifiable without opening the title page. Empty by default.
% -----------------------------------------------------------------------------
\makeatletter
\newcommand{\coursecode}[1]{\gdef\@coursecodetext{#1}}
\gdef\@coursecodetext{}
\makeatother

% -----------------------------------------------------------------------------
% Beamer theme assignments (only apply under Beamer).
%
% We detect Beamer via \@ifclassloaded{beamer}, which is the documented,
% non-internal way. This must be wrapped in \makeatletter/\makeatother
% because \@ifclassloaded contains an @-catcode character.
% -----------------------------------------------------------------------------
\makeatletter
\@ifclassloaded{beamer}{%
  \usefonttheme{professionalfonts}
  \setbeamercolor{structure}{fg=primary-blue}
  \setbeamercolor{frametitle}{fg=primary-blue}
  \setbeamercolor{title}{fg=primary-blue}
  \setbeamercolor{subtitle}{fg=primary-gold}
  \setbeamercolor{author}{fg=primary-blue}
  \setbeamercolor{institute}{fg=neutral}
  \setbeamercolor{date}{fg=primary-gold}
  \setbeamercolor{itemize item}{fg=primary-blue}
  \setbeamercolor{itemize subitem}{fg=primary-gold}
  \setbeamercolor{alerted text}{fg=primary-gold}
  \setbeamercolor{block title}{fg=white, bg=primary-blue}
  \setbeamercolor{block body}{bg=light-bg}
  % Semantic block variants: block=definition/concept (blue), exampleblock=
  % worked example (green/positive), alertblock=key takeaway or caution (gold).
  % Keep to <=2 colored boxes per slide across ALL three variants (INV-7).
  \setbeamercolor{block title example}{fg=white, bg=positive}
  \setbeamercolor{block body example}{bg=light-bg}
  \setbeamercolor{block title alerted}{fg=white, bg=primary-gold}
  \setbeamercolor{block body alerted}{bg=light-bg}

  % Minimal footer: course code (left) + page X / N (right), no navigation clutter
  \setbeamertemplate{navigation symbols}{}
  \setbeamertemplate{footline}{%
    \color{neutral}\footnotesize\hspace{0.7em}\@coursecodetext\hfill
    \insertframenumber/\inserttotalframenumber\hspace{0.7em}\vspace{0.3em}}
}{}
\makeatother

% -----------------------------------------------------------------------------
% TikZ setup — libraries the snippet gallery uses
% -----------------------------------------------------------------------------
\usepackage{tikz}
\usetikzlibrary{arrows.meta, positioning, calc, decorations.pathreplacing,
                fit, shapes.geometric, backgrounds}

% Shared TikZ styles — snippets and hand-written diagrams can pick these up
% via `\begin{tikzpicture}[every node/.style={...}]` or by naming specific
% styles (e.g., `\node[dag-node] (X) at (0,0) {X};`).
\tikzset{
  % Node styles for causal DAGs and similar
  dag-node/.style = {
    draw=primary-blue, fill=light-bg, rounded corners=2pt,
    minimum width=1.2cm, minimum height=0.9cm, align=center, font=\small
  },
  decision-node/.style = {
    draw=primary-gold, fill=white, diamond, aspect=1.8,
    minimum width=1.8cm, minimum height=1.2cm, align=center, font=\small,
    inner sep=1pt
  },
  % Edge styles — observed vs counterfactual
  observed-edge/.style = {-{Latex[length=2.5mm]}, thick, draw=jet},
  counterfactual-edge/.style = {-{Latex[length=2.5mm]}, thick, draw=neutral, dashed},
  confound-edge/.style = {-{Latex[length=2.5mm]}, thick, draw=negative, dashed},
  % Data-point styles for plots
  observed-dot/.style = {fill=primary-blue, draw=primary-blue, circle, inner sep=1.2pt},
  counterfactual-dot/.style = {fill=white, draw=neutral, circle, inner sep=1.2pt}
}

% -----------------------------------------------------------------------------
% Convenience macros
% -----------------------------------------------------------------------------
\newcommand{\muted}[1]{\textcolor{neutral}{#1}}
\newcommand{\key}[1]{\textcolor{primary-gold}{\textbf{#1}}}
\newcommand{\good}[1]{\textcolor{positive}{#1}}
\newcommand{\bad}[1]{\textcolor{negative}{#1}}

% Standout frame macro: full-bleed dark background for section transitions.
% Beamer-only (uses \begin{frame}/\setbeamercolor/\usebeamerfont) -- do not
% call from an article-class file (e.g. Notes/<CODE>/*-notes.tex). \muted,
% \key, \good, \bad, and \coursecode above are plain xcolor/gdef and are
% safe to use from either class; verified when Notes/article-class support
% was added.
% Expand: \transitionslide{Your transition title}
\newcommand{\transitionslide}[1]{%
  \begingroup
  \setbeamercolor{background canvas}{bg=primary-blue}
  \begin{frame}[plain]
    \centering
    \vfill
    {\usebeamerfont{title}\color{white}\Large #1\par}
    \vfill
  \end{frame}
  \endgroup
}

% Section divider frame (Beamer-only), an alternative to \transitionslide with a
% darker two-line style: near-black (jet) background, a small white label line
% above a larger gold title, and a thin gold rule along the bottom edge.
% Expand: \sectiondivider{Section 2}{Pointers}
\newcommand{\sectiondivider}[2]{%
  \begingroup
  \setbeamercolor{background canvas}{bg=jet}
  \begin{frame}[plain]
    \vspace*{3.0cm}
    {\color{white}\ttfamily\large #1\par}
    \vspace{0.5cm}
    {\color{primary-gold}\LARGE #2\par}
    \begin{tikzpicture}[remember picture, overlay]
      \draw[primary-gold, line width=2.5pt]
        ($(current page.south west)+(0,0.1)$) -- ($(current page.south east)+(0,0.1)$);
    \end{tikzpicture}
  \end{frame}
  \endgroup
}

% -----------------------------------------------------------------------------
% Channel watermark — "Concepts That Click" (YouTube).
%
% Article-class files (CompetitiveExam Q/A sets, Notes, Assignments) get a
% light diagonal draft watermark on every page plus a centered footer naming
% the channel. Beamer decks get a subtle TikZ background watermark instead
% (the footline already carries the course code). Centralized here so every
% MCQ PDF is branded without per-file edits.
% -----------------------------------------------------------------------------
\makeatletter
\@ifclassloaded{article}{%
  \usepackage{draftwatermark}
  \SetWatermarkText{Concepts That Click}
  \SetWatermarkScale{1}
  \SetWatermarkFontSize{2cm}
  \SetWatermarkLightness{0.86}
  \SetWatermarkAngle{45}
  \usepackage{fancyhdr}
  \pagestyle{fancy}
  \fancyhf{}
  \fancyfoot[C]{\footnotesize\textcolor{neutral}{Concepts That Click~|~YouTube: \href{https://www.youtube.com/@concepts.thatclick}{@concepts.thatclick}}}
  \renewcommand{\headrulewidth}{0pt}
  \renewcommand{\footrulewidth}{0pt}
  \fancypagestyle{plain}{%
    \fancyhf{}%
    \fancyfoot[C]{\footnotesize\textcolor{neutral}{Concepts That Click~|~YouTube: \href{https://www.youtube.com/@concepts.thatclick}{@concepts.thatclick}}}%
    \renewcommand{\headrulewidth}{0pt}%
    \renewcommand{\footrulewidth}{0pt}%
  }
}{}
\@ifclassloaded{beamer}{%
  \setbeamertemplate{background}{%
    \begin{tikzpicture}[remember picture, overlay]
      \node[rotate=45, opacity=0.055, align=center]
        at (current page.center)
        {\Huge\textcolor{neutral}{Concepts That Click}};
    \end{tikzpicture}%
  }
}{}
\makeatother 
\coursecode{JKSSB}
\title{Lesson 53 Practice: Networking Legacy Services and Web Standards}
\author{JKSSB Computer Awareness}
\date{\today}
\begin{document}
\maketitle
\noindent\textbf{Provenance:} All 20 questions are original JKSSB-pattern
questions (\textit{[Original]}); none reproduces an official paper item.

\begin{enumerate}[label=\textbf{Q\arabic*.}, leftmargin=*]
\item \textit{[Original]} Which switching method reserves a dedicated path
  for the whole session before data moves?
  \begin{enumerate}[label=\Alph*.]
    \item Packet switching
    \item Circuit switching
    \item Message flooding
    \item Address switching
  \end{enumerate}
\item \textit{[Original]} Which switching method breaks data into packets
  that travel independently, with no path reserved in advance?
  \begin{enumerate}[label=\Alph*.]
    \item Circuit switching
    \item Leased-line switching
    \item Packet switching
    \item Port switching
  \end{enumerate}
\item \textit{[Original]} Expand Telnet in its networking sense.
  \begin{enumerate}[label=\Alph*.]
    \item Telephone Network Terminal
    \item TELecommunication NETwork
    \item Television Network Transfer
    \item Terminal Network Telephone
  \end{enumerate}
\item \textit{[Original]} Which service provides remote login to another
  computer so the user can run commands on it?
  \begin{enumerate}[label=\Alph*.]
    \item FTP
    \item SMTP
    \item Telnet
    \item DNS
  \end{enumerate}
\item \textit{[Original]} Expand FTP.
  \begin{enumerate}[label=\Alph*.]
    \item Fast Text Processing
    \item File Transfer Protocol
    \item Full Terminal Procedure
    \item Forward Transmission Path
  \end{enumerate}
\item \textit{[Original]} Which service is used to upload and download files
  between computers on a network?
  \begin{enumerate}[label=\Alph*.]
    \item Telnet
    \item FTP
    \item SMTP
    \item IMAP
  \end{enumerate}
\item \textit{[Original]} Match the mail protocol to its role: which protocol
  \textbf{sends} mail from the sender toward the receiver's server?
  \begin{enumerate}[label=\Alph*.]
    \item SMTP
    \item POP3
    \item IMAP
    \item FTP
  \end{enumerate}
\item \textit{[Original]} Which protocol downloads mail to the local machine,
  typically removing it from the server?
  \begin{enumerate}[label=\Alph*.]
    \item SMTP
    \item IMAP
    \item POP3
    \item Telnet
  \end{enumerate}
\item \textit{[Original]} Which protocol keeps mail on the server and syncs
  it across several devices?
  \begin{enumerate}[label=\Alph*.]
    \item POP3
    \item SMTP
    \item FTP
    \item IMAP
  \end{enumerate}
\item \textit{[Original]} In an e-mail address, the symbol @ separates:
  \begin{enumerate}[label=\Alph*.]
    \item domain name and host name
    \item user name and domain name
    \item sender name and receiver name
    \item protocol name and address
  \end{enumerate}
\item \textit{[Original]} Expand W3C.
  \begin{enumerate}[label=\Alph*.]
    \item World Wide Web Consortium
    \item World Web Communication Centre
    \item Wide World Web Committee
    \item Web World Control Council
  \end{enumerate}
\item \textit{[Original]} Which of the following correctly describes DHTML?
  \begin{enumerate}[label=\Alph*.]
    \item A separate replacement language that replaces HTML entirely
    \item HTML combined with styling and scripting for dynamic pages
    \item A protocol for transferring files between servers
    \item A hardware socket for network cables
  \end{enumerate}
\item \textit{[Original]} Which of the following is \textbf{NOT} a correct
  statement about this lesson's services?
  \begin{enumerate}[label=\Alph*.]
    \item Telnet provides remote login to another computer.
    \item FTP transfers files between computers.
    \item SMTP sends mail toward the receiver's server.
    \item POP3 keeps mail synced on the server across devices.
  \end{enumerate}
\item \textit{[Original]} Which of the following statements is
  \textbf{INCORRECT}?
  \begin{enumerate}[label=\Alph*.]
    \item IPv6 addresses are 128 bits long.
    \item A digital signature proves authenticity and integrity of a message.
    \item ICMP is used to transfer application data between hosts.
    \item BGP routes traffic between autonomous systems.
  \end{enumerate}
\item \textit{[Original]} How many layers does the OSI reference model have?
  \begin{enumerate}[label=\Alph*.]
    \item Four
    \item Five
    \item Six
    \item Seven
  \end{enumerate}
\item \textit{[Original]} Which pairing is correctly matched?
  \begin{enumerate}[label=\Alph*.]
    \item Physical port --- a number identifying a network service
    \item Logical port --- a hardware socket for cables
    \item Physical port --- a hardware socket for cables or devices
    \item URL --- the numeric address of a device
  \end{enumerate}
\item \textit{[Original]} IPv4 addresses are 32 bits long, whereas IPv6
  addresses are \underline{\hspace{2.2cm}} long.
  \begin{enumerate}[label=\Alph*.]
    \item 64-bit
    \item 128-bit
    \item 256-bit
    \item 512-bit
  \end{enumerate}
\item \textit{[Original]} Which sequence correctly orders the chain from a
  typed web address to the reached machine?
  \begin{enumerate}[label=\Alph*.]
    \item IP address, DNS, URL
    \item URL typed, DNS resolves the domain name, browser reaches the IP address
    \item DNS typed, URL resolves the IP address, browser reaches the domain name
    \item URL reaches the IP address, which resolves DNS
  \end{enumerate}
\item \textit{[Original]} Which statement correctly distinguishes a physical
  port from a logical port?
  \begin{enumerate}[label=\Alph*.]
    \item A physical port is a service number; a logical port is a hardware socket.
    \item Both are hardware sockets for cables.
    \item A physical port is a hardware socket; a logical port is a number
      identifying a service or process.
    \item Both are numbers identifying a service or process.
  \end{enumerate}
\item \textit{[Original]} Evaluate the statements:
  \begin{enumerate}[label=\Roman*.]
    \item POP3 downloads mail to the local machine.
    \item A digital signature encrypts the message body so only the recipient can read it.
    \item HTTP/3 runs over QUIC over UDP, not over TCP.
  \end{enumerate}
  \begin{enumerate}[label=\Alph*.]
    \item I only
    \item I and III only
    \item II and III only
    \item I, II, and III
  \end{enumerate}
\end{enumerate}
\end{document}
